\section{Session Findings}

\subsection{Session 1}
\label{apx-s1}
\begin{enumerate}
\item \textbf{Biophilic experience was described as multi-channel and variable, in contrast to reduced sensory breadth indoors.} Experts described outdoor experience as involving concurrent sensory input—movement, texture, sound, smell, airflow, temperature, and depth of view—while noting that such concurrent sensory variation was reduced or absent indoors.

\item \textbf{Indoor environments were described as perceptually flattened and stabilised.}  Contemporary interiors were characterised as smooth, regular, and predictable, with limited variation across surfaces, routes, and environmental conditions.

\item \textbf{Bodily engagement present outdoors was described as diminished indoors.}  Experts noted that changes in gait, balance, and bodily effort commonly encountered outdoors were largely missing from indoor circulation and movement.

\item \textbf{Underfoot texture was identified as a missing experiential quality indoors.}  Fine-grained variation in floor texture, commonly encountered outdoors, was described as rarely present in indoor environments.

\item \textbf{View distance and layered visual motion were described as reduced indoors.}  Experts noted that extended views and subtle visual movement, such as shifting light or layered depth, were less available in indoor settings.

\item \textbf{Smell and air movement were described as absent or muted indoors.}  Freshness, seasonal smells, and airflow dynamics were repeatedly named as experiential cues present outdoors but missing indoors.

\item \textbf{Irregularity was tolerated outdoors but constrained indoors.}  Experts described unevenness and messiness as acceptable outdoors, while noting that similar qualities indoors were only tolerated when clearly intentional.

\item \textbf{Green elements alone were described as insufficient to address missing experiential qualities.}  Experts stressed that plants or green elements, when introduced without accompanying sensory and bodily variation, did not address the experiential absences they identified.

\item \textbf{Sound was described as situational rather than pervasive indoors.}  Experts differentiated between contexts where sound was useful (e.g., orientation or transition) and those where it was considered inappropriate as a continuous layer.
\end{enumerate}



\subsection{Session 2}
\label{apx-s2}
\begin{enumerate}
\item \textbf{Using vertical and transitional spaces to stage environmental change through sound.}  Stairs, corridors, and entry zones were structured so that shifts in acoustic character occur at moments of ascent, descent, or passage. Sound—alongside light and smell—was used to translate outdoor experiences of crossing into a new environment into indoor movement, rather than presenting a single static condition.

\item \textbf{Using soft material layering to create spatial depth.}  Curtains, screens, and other semi-transparent fabric layers were introduced as successive material layers that light and air pass through in sequence, translating outdoor depth and partial occlusion into interior space.

\item \textbf{Using filtered daylight to translate outdoor light variation.}  Daylight was filtered and reflected through architectural openings and surfaces to reproduce outdoor light shimmer and variation, avoiding focal displays and maintaining compatibility with everyday indoor use.

\item \textbf{Using gradual environmental drift to translate outdoor temporality.}  Slow changes in light, airflow, and acoustic character were embedded across daily or seasonal cycles, translating outdoor temporal variation into indoor environments without requiring explicit interaction or event-based control.

\item \textbf{Using exposed and operable building elements to express environmental responsiveness through sound.}  Façades, ducts, and other operable building systems were deliberately revealed so that changes in airflow, pressure, or use generate audible cues, allowing environmental responsiveness to be perceived through sound rather than interfaces.

\item \textbf{Using organised spatial irregularity to translate natural randomness.}  Experts translated outdoor experiences of natural randomness into indoor layouts through softly structured irregularities, such as multiple minor routes, offset elements, and subtle spatial deviations. These arrangements were designed to feel emergent rather than over-programmed, while remaining legible and navigable.

\item \textbf{Using spatial contrast and gradients to support state-dependent occupation.}  Experts organised spaces to offer parallel conditions—such as softer versus crisper textures, quieter versus livelier zones, or enclosed versus open areas—allowing occupants to position themselves relative to their current state without requiring explicit control.
\end{enumerate}

\subsection{Session 3}
\label{apx-s3}
\begin{enumerate}
\item \textbf{Sound was explored as a spatially distributed environmental field rather than a discrete source.}  
Experts consistently imagined biophilic sound as layered acoustic fields that extend across space, allowing occupants to enter, pass through, linger within, or leave them. Sound was thus framed as something to inhabit rather than attend to or control directly.

\item \textbf{Sound was used to structure space through depth, distance, and enclosure.}  
Differences in reverberation, diffusion, and acoustic scale were used to distinguish open from enclosed conditions and foreground from background activity, translating spatial depth and enclosure into sound.

\item \textbf{Sound was explored as a medium for conveying information and supporting habits and routines.}  
Gradual changes in sound density, intensity, and pattern across the day were developed to communicate information such as time, activity, and environmental state, enabling occupants to learn and interpret indoor soundscapes over time.

\item \textbf{Interaction with sound was framed as movement, proximity, and dwelling rather than control.}  
Rather than adjusting settings or mixing audio, occupants were imagined to engage with sound through bodily movement—approaching or retreating from sound-producing elements, choosing paths through differently tuned zones, or pausing within quieter acoustic pockets.

\item \textbf{Sound production and modulation were grounded in material contact and environmental processes.}  
Experts favoured sound that emerges through interactions between materials and conditions, such as water touching stone, wind moving foliage, or footsteps on softened surfaces, allowing sound to respond to environmental change rather than rely on abstract playback.

\item \textbf{Quiet was reframed as layered and natural rather than silent.}  
Instead of pursuing silence or full acoustic isolation, experts explored soundscapes that soften echo and murmur while preserving subtle biophilic layers, positioning natural quietness as an alternative to the acoustic vacuum of sealed interiors.
\end{enumerate}